\documentclass[11pt,a4paper]{article}
\usepackage[italian]{babel}
\usepackage[utf8]{inputenc}
\usepackage[T1]{fontenc}
\usepackage[margin=2.5cm]{geometry}
\usepackage{hyperref}
\usepackage{listings}
\usepackage{xcolor}
\usepackage{enumitem}
\usepackage{graphicx}
\let\oldunderscore\_
\renewcommand{\_}{\oldunderscore\allowbreak}

\lstset{
  basicstyle=\ttfamily\small,
  breaklines=true,
  frame=single,
  backgroundcolor=\color{gray!10},
  columns=fullflexible,
  keepspaces=true,
  upquote=true,
  showstringspaces=false,
}

\title{Guida passo-passo: robot a due ruote in equilibrio\\\large Pipeline standard NVIDIA (FreeCAD $\to$ URDF $\to$ Isaac Lab $\to$ RL)}
\author{Progetto robot\_slalom}
\date{Aggiornata al 24 ottobre 2026}

\begin{document}
\maketitle
\sloppy

\section*{Scopo di questa guida}
Documenta, passo per passo e in ordine cronologico, la pipeline standard NVIDIA/Isaac Lab usata in questo
progetto: dal modello CAD fino all'addestramento con reinforcement learning. Viene aggiornata ad ogni
passo compiuto. La regola generale della pipeline (cosa \`e "standard" e cosa no) \`e in
\texttt{PIPELINE\_STANDARD\_NVIDIA.md}.

\tableofcontents

\section{Punto di partenza: dal CAD Fusion a FreeCAD}
Il modello del robot a due ruote \`e stato progettato in \textbf{Fusion 360} (CAD): corpo, motori, ruote,
batteria, elettronica (ESP32-S3, sensori ToF, camera). Da qui si parte seguendo alla lettera la pipeline
standard NVIDIA, usando \textbf{FreeCAD} al posto di Fusion 360 per il passaggio successivo, perch\'e ha un
esportatore URDF nativo tramite il workbench \textbf{CROSS}:

\begin{enumerate}[leftmargin=*]
  \item \textbf{Importazione}: il modello viene esportato da Fusion in formato neutro (\texttt{STEP},
    file \texttt{export\textbackslash robot\_2ruote\_esp32\_varD.step}) e importato in FreeCAD.
  \item \textbf{CROSS}: dentro FreeCAD si definiscono i link e i giunti del robot (corpo, ruota sx, ruota
    dx, giunti rivoluti attorno all'asse delle ruote) e si assegnano materiali/densit\`a ai pezzi.
  \item \textbf{Esportazione}: CROSS esporta direttamente \texttt{robot.urdf}, con masse, inerzie e mesh
    incluse.
  \item \textbf{Simulazione}: Isaac Sim importa l'URDF (\texttt{sim\_utils.UrdfFileCfg}) e lo converte in
    USD.
  \item \textbf{Isaac Lab}: \texttt{ArticulationCfg} carica il robot, \texttt{DirectRLEnv} definisce
    l'ambiente RL, \texttt{rsl\_rl} esegue l'addestramento PPO.
\end{enumerate}

Tutto il codice Python di supporto (configurazione del robot, ambiente RL, script di training,
visualizzazione e valutazione) verr\`a scritto da zero a partire dal nuovo URDF, seguendo solo l'API
standard di Isaac Lab, senza bridge o file intermedi scritti a mano.

\section{Installazione FreeCAD e del server MCP}
\textbf{Data: 24 ottobre 2026.}

\begin{enumerate}[leftmargin=*]
  \item FreeCAD 1.1 installato in \texttt{C:\textbackslash Program Files\textbackslash FreeCAD 1.1}.
  \item Server MCP \texttt{freecad-robust-mcp} (progetto \href{https://github.com/spkane/freecad-addon-robust-mcp-server}{spkane/freecad-addon-robust-mcp-server})
    installato con pip (Python 3.11, richiesto per compatibilit\`a ABI con FreeCAD):
\begin{lstlisting}[language=bash]
python -m pip install --user freecad-robust-mcp
\end{lstlisting}
  \item Registrato in Claude Code:
\begin{lstlisting}[language=bash]
claude mcp add freecad "C:\Users\aless\AppData\Roaming\Python\Python311\Scripts\freecad-mcp.exe" -e FREECAD_MODE=xmlrpc
\end{lstlisting}
  \item Il server richiede la libreria \texttt{mcp} versione 1.x (con la 2.x si chiude all'avvio):
\begin{lstlisting}[language=bash]
python -m pip install --user "mcp<2"
\end{lstlisting}
  \item \textbf{Workbench dentro FreeCAD}: \emph{Robust MCP Bridge} non \`e nel catalogo dell'Addon
    Manager, quindi va installato a mano. Dal repository si copia la cartella
    \texttt{freecad\textbackslash RobustMCPBridge} in
\begin{lstlisting}
%APPDATA%\FreeCAD\v1-1\Mod\RobustMCPBridge\freecad\RobustMCPBridge\
\end{lstlisting}
    e si mette un \texttt{package.xml} in \texttt{...\textbackslash Mod\textbackslash RobustMCPBridge\textbackslash}.
    Attenzione: FreeCAD 1.1 legge gli addon dalla cartella con la versione (\texttt{v1-1\textbackslash Mod}),
    non da \texttt{FreeCAD\textbackslash Mod}.
  \item Il Python interno di FreeCAD deve avere le dipendenze del bridge (mancava \texttt{anyio}):
\begin{lstlisting}[language=bash]
"C:\Program Files\FreeCAD 1.1\bin\python.exe" -m pip install --user freecad-robust-mcp
\end{lstlisting}
  \item Avvio: in FreeCAD selezionare il workbench \emph{Robust MCP Bridge} e cliccare
    \emph{Start MCP Bridge}. Nel Report view compare ``MCP Bridge started!'' (XML-RPC 9875, socket 9876)
    e \texttt{claude mcp list} mostra \texttt{freecad: Connected}.
  \item Avvio automatico del bridge a ogni apertura di FreeCAD: in \texttt{\%APPDATA\%\textbackslash FreeCAD\textbackslash v1-1\textbackslash user.cfg},
    nel gruppo \texttt{RobustMCPBridge}, aggiungere \texttt{<FCBool Name="AutoStart" Value="1"/>}
    (con FreeCAD chiuso, altrimenti viene sovrascritto).
\end{enumerate}

\section{Installazione del workbench CROSS (export URDF)}
CROSS non \`e nel catalogo standard dell'Addon Manager (cercando ``cross'' escono solo \emph{Cross Section}
e \emph{Mouse Cross}). Si installa a mano dal repository \texttt{github.com/galou/freecad.cross}:
\begin{enumerate}[leftmargin=*]
  \item Clonare il repository nella cartella Mod, \textbf{con il nome esatto} \texttt{freecad.cross}
    (il codice cerca icone, template e finestre in \texttt{Mod\textbackslash freecad.cross\textbackslash resources}:
    con un altro nome il workbench si carica ma i template dell'export non vengono trovati):
\begin{lstlisting}[language=bash]
cd %APPDATA%\FreeCAD\v1-1\Mod
git clone --depth 1 https://github.com/galou/freecad.cross.git freecad.cross
\end{lstlisting}
  \item Creare \texttt{Mod\textbackslash freecad.cross\textbackslash InitGui.py} con una sola riga
    (l'Addon Manager lo farebbe da s\'e; col clone manuale va scritto):
\begin{lstlisting}[language=Python]
import freecad.cross.init_gui
\end{lstlisting}
  \item Installare le dipendenze dichiarate in \texttt{package.xml} nel Python di FreeCAD:
\begin{lstlisting}[language=bash]
"C:\Program Files\FreeCAD 1.1\bin\python.exe" -m pip install --user urdf-parser-py xacro
\end{lstlisting}
  \item Riavviare FreeCAD. Nell'elenco dei workbench compare \emph{CROSS - ROS workbench} con i comandi
    NewRobot, NewLink, NewJoint, SetMaterial, CalculateMassAndInertia, UrdfExport.
    L'avviso ``\texttt{ROS\_DISTRO} is not set'' \`e normale: ROS non serve per esportare l'URDF.
\end{enumerate}

\section{I concetti: Robot, Link, Joint}
Un robot per la simulazione (formato URDF) si descrive con tre soli mattoni, che in CROSS corrispondono
ai pulsanti \emph{New Robot}, \emph{New Link}, \emph{New Joint}:

\begin{description}[leftmargin=*]
  \item[Link = un pezzo rigido.] Un insieme di pezzi CAD che si muovono sempre insieme, come se fossero
    incollati. Per il simulatore \`e un unico oggetto con massa, baricentro e inerzia. Il link non ha una
    forma propria: \`e un'``etichetta'' che raggruppa pezzi CAD (per questo nell'albero di FreeCAD,
    passandoci sopra col mouse, non si evidenzia niente).
  \item[Joint = il collegamento tra due link.] Dice come un link si muove rispetto all'altro:
    \begin{itemize}
      \item \textbf{Parent}: il pezzo ``fisso''; \textbf{Child}: il pezzo che si muove.
      \item \textbf{Type}: \texttt{continuous} (gira all'infinito, come una ruota), \texttt{revolute}
        (gira entro dei limiti, come un gomito), \texttt{prismatic} (scorre, come un cassetto),
        \texttt{fixed} (bloccato).
      \item \textbf{Origin}: posizione del perno e direzione dell'asse di rotazione.
      \item \textbf{Effort / Velocity}: coppia (N$\cdot$m) e velocit\`a (rad/s) massime del motore.
    \end{itemize}
  \item[Robot = il contenitore] che tiene insieme link e giunti; \`e ci\`o che si esporta come
    \texttt{robot.urdf}.
\end{description}

In Isaac Lab la rete neurale comander\`a i \textbf{giunti} (quanta coppia dare a ogni ruota); il
simulatore calcola come si muove il resto in base alle masse dei \textbf{link}.

\section{Composizione del robot in CROSS}
\textbf{Data: 24 ottobre 2026.} File: \texttt{export\textbackslash robot\_cross.FCStd}
(recuperato dall'autosalvataggio di FreeCAD dopo uno spegnimento del PC, poi salvato con questo nome).

\subsection{Struttura}
Il robot \`e fatto di \textbf{3 link} e \textbf{2 joint}:
\begin{itemize}[leftmargin=*]
  \item \textbf{Link} (pezzi rigidi): \texttt{base\_link} (il telaio), \texttt{ruota\_sx} (ruota
    sinistra), \texttt{ruota\_dx} (ruota destra). \textbf{Anche ogni ruota \`e un link}: \`e un pezzo
    rigido a s\'e, separato dal telaio, perch\'e si muove (gira) rispetto ad esso.
  \item \textbf{Joint} (collegamenti): \texttt{giunto\_ruota\_sx} collega il telaio alla ruota sinistra,
    \texttt{giunto\_ruota\_dx} collega il telaio alla ruota destra.
\end{itemize}

\begin{lstlisting}
               [LINK] base_link  (telaio, 394 g)
                 /                       \
   [JOINT] giunto_ruota_sx       [JOINT] giunto_ruota_dx
        (continuous)                  (continuous)
              |                             |
      [LINK] ruota_sx               [LINK] ruota_dx
         (41 g)                        (41 g)
\end{lstlisting}
In generale: un link per ogni parte che si muove in modo indipendente, un joint tra ogni coppia di link
collegati. Ogni joint ha sempre un link \emph{parent} (sopra) e un link \emph{child} (sotto).

Regola usata: tutto ci\`o che \`e avvitato al telaio (anche i motori, il cui corpo non gira) va in
\texttt{base\_link}; solo ci\`o che ruota con l'albero del motore va nel link della ruota.

\subsection{Pezzi CAD in ogni link}
Lo STEP importato da Fusion produce 40 solidi (\texttt{robot\_2ruote\_esp32\_varD000}\ldots\texttt{039}),
senza i nomi dei corpi. Le etichette in FreeCAD sono i nomi dei corpi in Fusion: lo STEP mantiene lo
stesso ordine, e la corrispondenza \`e verificata pezzo per pezzo su dimensioni e volumi (es.
\texttt{esp32\_pcb} $69{,}4\times25{,}4\times1{,}6$ mm, i ToF $13\times25$ mm).

\begin{description}[leftmargin=*]
  \item[\texttt{base\_link}] (36 pezzi, 394,1 g, baricentro 3,6 mm avanti e 27,8 mm sopra l'asse delle
    ruote, cio\`e $z = 60{,}3$ mm):
    \begin{itemize}
      \item Struttura: piastra\_alta, piastra\_bassa, colonnina\_1\ldots4, portabatterie,
        appoggio\_anteriore, appoggio\_posteriore.
      \item Motori e staffe: motore\_sx, motore\_dx, staffa\_alta\_sx/dx, staffa\_base\_sx/dx.
      \item Batteria: cella\_18650\_1/2.
      \item Elettronica: driver\_TB6612, step\_down\_5V, IMU\_MPU6050, esp32\_pcb, esp32\_modulo\_WROOM,
        esp32\_usb\_1/2, esp32\_header\_sx/dx, esp32\_supporto\_1/2.
      \item Sensori di distanza: ToF\_centro, ToF\_sx, ToF\_dx.
      \item Camera: cam\_supporto, cam\_pcb, cam\_modulo\_WROOM, cam\_obiettivo\_base, cam\_obiettivo.
    \end{itemize}
  \item[\texttt{ruota\_sx}] (41,2 g, centro $y{=}{+}88$, $z{=}32{,}5$ mm): ruota\_sx\_corpo (cerchio
    stampato), ruota\_sx\_gomma (camera d'aria).
  \item[\texttt{ruota\_dx}] (41,2 g, centro $y{=}{-}88$, $z{=}32{,}5$ mm): ruota\_dx\_corpo,
    ruota\_dx\_gomma.
\end{description}

\subsection{Giunti}
\begin{center}
\resizebox{\textwidth}{!}{%
\begin{tabular}{lllll}
  \hline
  Giunto & Tipo & Parent $\to$ Child & Origin (mm, rotazione) & Effort / Velocity \\
  \hline
  \texttt{giunto\_ruota\_sx} & continuous & base\_link $\to$ ruota\_sx & (0, +88, 32,5), roll $-90^\circ$ & 0,30 N$\cdot$m / 29,3 rad/s \\
  \texttt{giunto\_ruota\_dx} & continuous & base\_link $\to$ ruota\_dx & (0, $-$88, 32,5), roll $-90^\circ$ & 0,30 N$\cdot$m / 29,3 rad/s \\
  \hline
\end{tabular}}
\end{center}
La rotazione di $-90^\circ$ porta l'asse $z$ del giunto (asse di rotazione in URDF) lungo l'asse delle
ruote. Effort e Velocity sono i valori \textbf{tipici, non misurati} del motore JGA25-370 a 6 V:
coppia di stallo 0,30 N$\cdot$m e 280 giri/min a vuoto ($280 \cdot 2\pi/60 = 29{,}3$ rad/s).

\subsection{Come si crea dall'editor}
\begin{enumerate}[leftmargin=*]
  \item Workbench \emph{CROSS - ROS workbench}.
  \item \emph{New Robot} (\texttt{N},\texttt{R}): crea il contenitore \texttt{robot}.
  \item \emph{New Link} (\texttt{N},\texttt{L}) con il robot selezionato; rinominare (F2) in
    \texttt{base\_link}. Nel pannello \emph{Dati}, propriet\`a \textbf{Real} $\to$ ``\ldots'' e spuntare i
    pezzi elencati sopra; \textbf{Visual} e \textbf{Collision} con gli stessi pezzi.
  \item Stessa cosa per \texttt{ruota\_sx} e \texttt{ruota\_dx}.
  \item \emph{New Joint} (\texttt{N},\texttt{J}): nel pannello Dati Type, Parent, Child, Origin,
    Effort, Velocity come in tabella. Usare \emph{Set placement} per posizionare le mesh.
  \item Massa e inerzia di ogni link: con lo script della sezione \ref{sec:masse} nella console Python di
    FreeCAD (non con il pulsante \emph{Calculate mass and inertia}, vedi l\`i perch\'e).
\end{enumerate}

\subsection{Problemi incontrati}
\begin{itemize}[leftmargin=*]
  \item \textbf{Si vedono solo frecce}: sui link \texttt{ShowVisual} e \texttt{ShowReal} (pannello
    \emph{Vista}) erano \texttt{False}. Messi a \texttt{True}, e nascosti i 40 pezzi CAD originali
    per non vederli doppi.
  \item \textbf{Ruote staccate}: i pezzi CAD sono in coordinate assolute, ma il link applicava di nuovo la
    propria posizione ($\pm$88 mm, $-90^\circ$). Corretto impostando su \texttt{ruota\_sx/dx}
    \texttt{MountedPlacement} = inverso del \texttt{Placement} del link (\`e ci\`o che fa
    \emph{Set placement}).
  \item \textbf{Colori} per riconoscere i link: grigio \texttt{base\_link}, rosso \texttt{ruota\_sx}, blu
    \texttt{ruota\_dx}. Le frecce blu al centro delle ruote sono gli assi dei giunti
    (\texttt{ShowAxis}, lunghezza 60 mm).
  \item \textbf{Nomi dei pezzi sbagliati}: nella prima sessione i 40 solidi erano stati rinominati a
    intuito (es.\ i 3 sensori ToF chiamati \texttt{supporto\_anteriore\_*}, driver/step-down/IMU chiamati
    \texttt{scheda\_piccola\_*} e \texttt{componente\_1}). Rinominati con i nomi dei corpi Fusion.
  \item \textbf{Masse di origine non documentata}: masse e inerzie erano state calcolate con densit\`a
    non scritte da nessuna parte (393 g + 2$\times$42 g). Ricalcolate con lo script della sezione
    \ref{sec:masse} dalle masse per pezzo stimate in Fusion: ora coincidono con i valori di Fusion.
  \item \textbf{Baricentro e inerzia delle ruote sbagliati nell'URDF}: il primo export dava per
    \texttt{ruota\_sx/dx} il baricentro a $(0;\pm0{,}088;0{,}0325)$ m \emph{rispetto alla ruota} (cio\`e
    in coordinate assolute) e l'inerzia maggiore su \texttt{iyy}, mentre nel sistema del link ruota
    l'asse di rotazione \`e $z$. Lo script della sezione \ref{sec:masse} esprime baricentro e inerzia nel
    sistema di ogni link.
  \item \textbf{Export: la finestra ``ROS workspace''}: all'avvio di \emph{Export to URDF} CROSS chiede
    il workspace ROS (non la cartella di uscita!) e avvisa ``\ldots does not appear to be a valid ROS
    workspace''. \`E innocuo: controlla solo \texttt{install/setup.bash} di ROS~2 su Linux. Scegliere
    \texttt{G:\textbackslash robotTest}. La cartella di uscita \`e solo l'\textbf{OutputPath} del robot.
  \item \textbf{Export che non scrive nulla}: se i file esistono gi\`a, nella finestra di conferma vanno
    messi tutti su \textbf{Overwrite}; con \emph{Ignore} CROSS non scrive niente (e non accetta una
    scrittura parziale). Finch\'e quella finestra resta aperta, FreeCAD (e il bridge MCP) \`e bloccato.
  \item \textbf{Bridge MCP che non parte da solo}: all'avvio compare
    ``\texttt{module 'FreeCAD' has no attribute 'GuiUp'}'' e l'auto-start fallisce; avviarlo dal
    workbench \emph{Robust MCP Bridge} con \emph{Start MCP Bridge}.
\end{itemize}

\section{Calcolo di masse e inerzie}\label{sec:masse}
\subsection{Perch\'e non il pulsante di CROSS}
\emph{Calculate mass and inertia} di CROSS (verificato nel codice) usa \textbf{solo il primo pezzo} della
lista Real di ogni link e \textbf{una sola densit\`a per link}, e lascia l'inerzia negli assi del CAD. Per
\texttt{base\_link} (36 pezzi di materiali diversi) e per le ruote (cerchio + gomma) il risultato sarebbe
sbagliato. Per questo i link hanno \texttt{MaterialNotCalculate = True}: il pulsante li salta.

\subsection{Masse dei pezzi}
Il dato di partenza \`e la massa di ogni pezzo. Per ora sono stime fatte in Fusion; \textbf{quando ci sono i
pezzi veri si pesano e si correggono questi numeri}.
\begin{center}
\resizebox{\textwidth}{!}{%
\begin{tabular}{lrl}
  \hline
  Pezzo & Massa (g) & Origine della stima \\
  \hline
  piastra\_alta, piastra\_bassa & 32,60; 32,63 & volume $\times$ 0,6 g/cm$^3$ (PLA stampato) \\
  colonnina\_1\ldots4 & 0,74 cad. & volume $\times$ 0,6 g/cm$^3$ \\
  staffa\_alta/base\_sx/dx & 7,39 cad. & volume $\times$ 0,6 g/cm$^3$ \\
  portabatterie & 9,34 & volume $\times$ 0,6 g/cm$^3$ \\
  appoggio\_anteriore / posteriore & 1,27; 0,97 & volume $\times$ 0,6 g/cm$^3$ \\
  motore\_sx, motore\_dx & 70 cad. & motore equivalente JGA25-370 (\textbf{da pesare}) \\
  cella\_18650\_1/2 & 46 cad. & datasheet (dato reale) \\
  driver\_TB6612 / step\_down\_5V / IMU\_MPU6050 & 5 / 10 / 3 & stima \\
  ToF\_centro/sx/dx & 3 cad. & stima \\
  esp32\_pcb / esp32\_modulo\_WROOM & 8 / 3 & stima \\
  esp32\_usb\_1/2, esp32\_header\_sx/dx & 0,5 / 1 cad. & stima \\
  esp32\_supporto\_1/2 & 0,13 cad. & volume $\times$ 0,6 g/cm$^3$ \\
  cam\_supporto / cam\_pcb / cam\_modulo\_WROOM & 0,54 / 5 / 3 & stima \\
  cam\_obiettivo\_base / cam\_obiettivo & 1 / 2 & stima \\
  ruota\_sx/dx\_corpo & 37,06 cad. & volume $\times$ 0,6 g/cm$^3$ \\
  ruota\_sx/dx\_gomma & 4,11 cad. & volume $\times$ 0,95 g/cm$^3$ (gomma) \\
  \hline
  \textbf{Totale robot} & \textbf{476,5} & base\_link 394,1 + 2 ruote $\times$ 41,2 \\
  \hline
\end{tabular}}
\end{center}

\subsection{Lo script}
In FreeCAD: \emph{Visualizza} $\to$ \emph{Pannelli} $\to$ \emph{Console Python}, incollare lo script e
premere Invio. Per ogni link: somma le masse dei pezzi, calcola il baricentro, somma le inerzie di ogni
pezzo (dalla sua forma, a densit\`a uniforme) portandole sul baricentro comune (teorema di
Huygens-Steiner), poi riporta baricentro e inerzia nel sistema del link (per le ruote: quello del giunto).
Infine scrive \texttt{Mass}, \texttt{CenterOfMass}, \texttt{Ixx}\ldots\texttt{Izz} nei link. Dopo, salvare e
rifare l'export URDF.
\begin{lstlisting}[language=Python,basicstyle=\ttfamily\scriptsize]
import FreeCAD

# Massa di ogni pezzo in grammi (stime da Fusion 360: da sostituire con i pesi veri).
MASSE_G = {
    "piastra_alta": 32.60, "piastra_bassa": 32.63,
    "colonnina_1": 0.74, "colonnina_2": 0.74, "colonnina_3": 0.74, "colonnina_4": 0.74,
    "staffa_alta_sx": 7.39, "staffa_base_sx": 7.39, "staffa_alta_dx": 7.39, "staffa_base_dx": 7.39,
    "portabatterie": 9.34, "appoggio_anteriore": 1.27, "appoggio_posteriore": 0.97,
    "motore_sx": 70.0, "motore_dx": 70.0,
    "cella_18650_1": 46.0, "cella_18650_2": 46.0,
    "driver_TB6612": 5.0, "step_down_5V": 10.0, "IMU_MPU6050": 3.0,
    "ToF_centro": 3.0, "ToF_sx": 3.0, "ToF_dx": 3.0,
    "esp32_supporto_1": 0.13, "esp32_supporto_2": 0.13, "esp32_pcb": 8.0, "esp32_modulo_WROOM": 3.0,
    "esp32_usb_1": 0.5, "esp32_usb_2": 0.5, "esp32_header_sx": 1.0, "esp32_header_dx": 1.0,
    "cam_supporto": 0.54, "cam_pcb": 5.0, "cam_modulo_WROOM": 3.0,
    "cam_obiettivo_base": 1.0, "cam_obiettivo": 2.0,
    "ruota_sx_corpo": 37.06, "ruota_dx_corpo": 37.06,
    "ruota_sx_gomma": 4.11, "ruota_dx_gomma": 4.11,
}

doc = FreeCAD.ActiveDocument
robot = [o for o in doc.Objects if getattr(o, "_Type", "") == "Cross::Robot"][0]
links = [o for o in robot.Group if getattr(o, "_Type", "") == "Cross::Link"]
joints = [o for o in robot.Group if getattr(o, "_Type", "") == "Cross::Joint"]

for link in links:
    # Sistema del link: origine del giunto di cui e' figlio (base_link: origine del robot).
    frame = FreeCAD.Placement()
    for j in joints:
        if j.Child == link.Name:
            frame = j.Origin

    m_tot, com = 0.0, FreeCAD.Vector()
    for p in link.Real:
        m = MASSE_G[p.Label] / 1000.0
        m_tot += m
        com += p.Shape.CenterOfMass * m
    com = com * (1.0 / m_tot)

    # Tensore d'inerzia attorno al baricentro del link, assi del CAD (kg*m^2).
    I = [[0.0] * 3 for _ in range(3)]
    for p in link.Real:
        m = MASSE_G[p.Label] / 1000.0
        k = m / p.Shape.Volume * 1e-6
        M = p.Shape.MatrixOfInertia
        Ip = [[M.A11, M.A12, M.A13], [M.A21, M.A22, M.A23], [M.A31, M.A32, M.A33]]
        r = (p.Shape.CenterOfMass - com) * 1e-3
        rv, r2 = [r.x, r.y, r.z], r.x**2 + r.y**2 + r.z**2
        for a in range(3):
            for b in range(3):
                I[a][b] += Ip[a][b] * k + m * ((r2 if a == b else 0.0) - rv[a] * rv[b])

    # Riporta baricentro e inerzia nel sistema del link.
    R = frame.Rotation.toMatrix()
    Rm = [[R.A11, R.A12, R.A13], [R.A21, R.A22, R.A23], [R.A31, R.A32, R.A33]]
    IL = [[sum(Rm[a][i] * I[a][b] * Rm[b][j] for a in range(3) for b in range(3))
           for j in range(3)] for i in range(3)]

    link.MaterialNotCalculate = True
    link.Mass = m_tot
    link.CenterOfMass = FreeCAD.Placement(frame.inverse().multVec(com), FreeCAD.Rotation())
    link.Ixx, link.Iyy, link.Izz = IL[0][0], IL[1][1], IL[2][2]
    link.Ixy, link.Ixz, link.Iyz = IL[0][1], IL[0][2], IL[1][2]
    print("%-10s massa %.1f g  baricentro %s mm  Ixx %.3e Iyy %.3e Izz %.3e"
          % (link.Label, m_tot * 1000, link.CenterOfMass.Base, link.Ixx, link.Iyy, link.Izz))

doc.recompute()
\end{lstlisting}

Risultato (uguale ai valori calcolati da Fusion 360 sul modello originale):
\begin{lstlisting}[basicstyle=\ttfamily\scriptsize]
base_link  massa 394.1 g  baricentro (3.6, -0.02, 60.3) mm  Ixx 9.033e-04 Iyy 5.547e-04 Izz 8.051e-04
ruota_sx   massa 41.2 g   baricentro (0, 0, 0) mm          Ixx 1.265e-05 Iyy 1.265e-05 Izz 2.257e-05
ruota_dx   massa 41.2 g   baricentro (0, 0, 0) mm          Ixx 1.265e-05 Iyy 1.265e-05 Izz 2.257e-05
\end{lstlisting}
Nota: lo script assume che il padre di ogni giunto sia \texttt{base\_link} con posizione zero (vero per
questo robot); con catene pi\`u lunghe andrebbero composte le posizioni dei giunti.

\section{Il formato URDF}
\textbf{URDF} (\emph{Unified Robot Description Format}) \`e un file di testo in formato XML che descrive
un robot come albero di \textbf{link} collegati da \textbf{joint}. Nato in ROS, \`e letto da quasi tutti
i simulatori, tra cui Isaac Sim. Il nostro \`e \texttt{robot\_description\textbackslash urdf\textbackslash robot.urdf}.

\textbf{Unit\`a di misura}: sempre SI. Lunghezze in metri, angoli in radianti, masse in kg, inerzie in
kg$\cdot$m$^2$ (CROSS converte dai millimetri di FreeCAD).

\subsection{Struttura del file}
\begin{lstlisting}
<robot name="robot">
  <link name="base_link"> ... </link>
  <link name="ruota_sx">  ... </link>
  <link name="ruota_dx">  ... </link>
  <joint name="giunto_ruota_sx" type="continuous"> ... </joint>
  <joint name="giunto_ruota_dx" type="continuous"> ... </joint>
</robot>
\end{lstlisting}

\subsection{Un link}
Ogni \texttt{<link>} ha tre parti (esempio reale, \texttt{ruota\_sx}, accorciato):
\begin{lstlisting}
<link name="ruota_sx">
  <visual>       <!-- come si VEDE: una per ogni pezzo CAD -->
    <origin rpy="1.5708 0 0" xyz="0 0.0325 -0.088"/>
    <geometry>
      <mesh filename="package://robot_description/meshes/robot_cross_ruota_sx_gomma.dae"/>
    </geometry>
  </visual>
  <collision>    <!-- forma usata per i CONTATTI (qui uguale alla visual) -->
    ...
  </collision>
  <inertial>     <!-- come si MUOVE: massa, baricentro, inerzia -->
    <mass value="0.0418"/>
    <origin xyz="0 0 0"/>
    <inertia ixx="1.30e-05" iyy="1.30e-05" izz="2.32e-05" ixy="0" ixz="0" iyz="0"/>
  </inertial>
</link>
\end{lstlisting}
\begin{itemize}[leftmargin=*]
  \item \texttt{<visual>}: solo grafica, non influisce sulla fisica.
  \item \texttt{<collision>}: la forma che tocca il pavimento e gli ostacoli. Pu\`o essere la mesh vera o una
    forma semplice (cilindro, scatola), pi\`u veloce da simulare.
  \item \texttt{<inertial>}: massa, posizione del baricentro (\texttt{origin}, nel sistema del link) e
    tensore d'inerzia attorno al baricentro. \`E ci\`o che conta per la dinamica.
  \item \texttt{<origin xyz rpy>}: spostamento (m) e rotazione (roll, pitch, yaw in radianti) rispetto al
    sistema di riferimento del link.
  \item \texttt{package://robot\_description/\ldots}: percorso della mesh secondo la convenzione ROS
    (nome del pacchetto = nome della cartella).
\end{itemize}

\subsection{Un joint}
\begin{lstlisting}
<joint name="giunto_ruota_sx" type="continuous">
  <parent link="base_link"/>
  <child  link="ruota_sx"/>
  <origin rpy="-1.5708 0 0" xyz="0 0.088 0.0325"/>
  <axis xyz="0 0 1"/>
</joint>
\end{lstlisting}
\begin{itemize}[leftmargin=*]
  \item \texttt{parent} / \texttt{child}: i due link collegati.
  \item \texttt{origin}: dove sta il perno, nel sistema del parent. Qui al centro della ruota sinistra
    (y = 88 mm, z = 32,5 mm), ruotato di $-90^\circ$ attorno a $x$.
  \item \texttt{axis}: asse di rotazione nel sistema del joint. $z$ dopo la rotazione di $-90^\circ$
    coincide con l'asse della ruota.
  \item \texttt{<limit effort velocity>}: coppia e velocit\`a massime. Per i giunti \texttt{continuous}
    sono facoltativi e CROSS non li scrive; in Isaac Lab si impostano negli attuatori
    dell'\texttt{ArticulationCfg}.
\end{itemize}

\textbf{Sistemi di riferimento}: il link radice (\texttt{base\_link}) definisce il sistema del robot;
ogni joint posiziona il sistema del link figlio rispetto al padre. Per questo baricentro e mesh delle ruote
sono espressi rispetto al centro della ruota, non rispetto all'origine del robot (vedi il problema
del baricentro delle ruote).

\section{Prossimi passi (da fare)}
\begin{enumerate}[leftmargin=*]
  \item Quando ci sono i pezzi veri: pesarli, correggere le masse e ricalcolare massa e inerzia.
  \item FATTO (24/10/2026): esportare l'URDF da CROSS: selezionare \texttt{robot} e premere \emph{Export to URDF},
    confermando la scrittura di tutti i file. La destinazione \`e la propriet\`a \textbf{OutputPath} del
    robot, impostata a \texttt{G:/robotTest/robot\_description}: una cartella dedicata solo all'output
    generato (sovrascritta a ogni export), separata da \texttt{robot\_sim\textbackslash} dove andr\`a il
    codice Isaac Lab scritto a mano. Il nome segue la convenzione ROS \texttt{<nome>\_description}; nell'URDF
    le mesh diventano \texttt{package://robot\_description/meshes/\ldots}. Servono solo \texttt{urdf\textbackslash}
    e \texttt{meshes\textbackslash}; \texttt{package.xml}, \texttt{CMakeLists.txt}, \texttt{launch\textbackslash},
    \texttt{rviz\textbackslash} sono file ROS non usati da Isaac.
  \item Scrivere da capo \texttt{robot\_cfg.py} (ArticulationCfg/UrdfFileCfg) e l'ambiente RL
    (DirectRLEnv) usando i nomi di link/giunti del nuovo URDF.
  \item Scrivere da capo gli script di training/visualizzazione/valutazione e confrontare i risultati con
    quelli precedenti (8,33\textdegree{} di appoggio, reward RL) per verificare l'equivalenza del nuovo
    modello.
\end{enumerate}

\end{document}
